\subsection{H3: Weighted Motif Feature Discrimination}\label{sec:results_h3}

Spectral clustering with weighted motif features (19 features) achieves $NMI=0.705$ ($ARI=0.516$, optimal $K=8$) against ground-truth domain labels, dramatically outperforming binary motif counts (4 features) at $NMI=0.101$ ($ARI=0.008$). This seven-fold NMI improvement is highly significant ($p=0.001$ by 1000-permutation test) and demonstrates that edge-weight information is essential for capability discrimination (Table~\ref{table:clustering}).

\paragraph{Systematic Feature Set Comparison.}
Weighted motif features ($NMI=0.705$) slightly outperform graph-level statistics ($NMI=0.677$, 8 features, optimal $K=4$); this difference, while modest in absolute terms, is statistically significant ($p=0.001$). Combining all features --- weighted motif, binary motif, and graph statistics (31 features) --- yields the best overall performance at $NMI=0.844$ ($ARI=0.688$, $K=8$), significantly exceeding the best single feature set ($p=0.001$). Notably, adding binary features to weighted features ($NMI=0.705$, 23 features) provides zero improvement over weighted alone, confirming that binary motif counts carry no information beyond what weighted features already capture. The improvement from adding graph statistics to weighted features (from 0.705 to 0.844) indicates that graph-level properties contribute complementary domain-relevant information.

\paragraph{Most Discriminative Individual Features.}
One-way ANOVA with domain as the grouping variable reveals that FFL path dominance ($\eta^2=0.917$) is the single most discriminative feature, explaining over 91\% of between-domain variance. FFL intensity mean ($\eta^2=0.810$) and FFL intensity Q25 ($\eta^2=0.835$) also show very large effect sizes, followed by Onnela coherence std ($\eta^2=0.747$). For comparison, node count achieves $\eta^2=0.927$, making it the single most discriminative graph-level feature. That a motif feature (path dominance) approaches node count in discriminative power is notable, given that node count directly reflects circuit complexity while path dominance captures a subtle structural property.

\paragraph{Honest Assessment of Partial Negative Results.}
While weighted features substantially improve over binary counts, several findings temper the conclusion. First, the gap between weighted-only and combined features (NMI 0.705 vs.\ 0.844) shows that motif features alone do not capture all domain-relevant structure. Second, graph-level statistics contribute complementary information, particularly node count and edge count which vary systematically across domains due to inherent differences in circuit complexity (e.g., multi-hop reasoning graphs are larger than antonym graphs). Third, the binary motif baseline is extremely weak ($NMI={gt['b_nmi']:.3f}$), suggesting that raw motif topology without weight information is nearly uninformative for domain discrimination in this corpus.

Our weighted motif feature framework adapts the intensity and coherence measures of \citet{onnela2005intensity} to signed, layered DAGs, extending classical weighted network analysis \cite{fagiolo2007clustering} to the specific structure of neural network attribution graphs.

\textbf{Verdict:} H3 is \textbf{supported with caveats}. Weighted motif features vastly outperform binary counts and slightly outperform graph statistics for domain clustering, but combined features are needed for best performance, indicating complementary information sources.